\section{The proof of Theorem 4}\label{s:proof}

The proof has three sentences. The first instantiates Assumption~3, the second infers that a life containing the resulting engagement consists in engagement to a positive extent, and the third applies Definition~1. The proof is valid given one principle that it uses without stating, and the comment on \Lm{7} identifies it.

\lem{6}{By Assumption 3, there is a project $p$ of objective worth in which the agent can engage.}

\paragraph{Instantiation.} The sentence performs two instantiations. It instantiates the universal quantifier of Assumption~3 to \enquote{the agent}, the arbitrary agent of Theorem~4's \enquote{every such agent}, and it introduces a name, $p$, for a project whose existence Assumption~3 then guarantees. The second is existential instantiation, and the proof treats $p$ as arbitrary among the projects that meet the condition: nothing is assumed about $p$ except that it is of objective worth and that the agent can engage in it. The Lean module performs the same step with \texttt{obtain} (Section~\ref{s:lean}).

\paragraph{\enquote{can} for \enquote{is able to}.} Assumption~3 says that the agent \enquote{is able to engage}; the proof says that she \enquote{can engage}. The substitution is harmless, but it marks that the modality of Assumption~3 is the modality of the ordinary \enquote{can} of ability, which the comment on \Lm{4.4} has analyzed.

\lem{7}{A life that includes engagement with $p$ consists, to some positive extent, in active engagement with a project of objective worth.}

This is the central step of the proof, and it contains more than its surface suggests. It is a universally quantified conditional: any life that includes engagement with $p$ consists to a positive extent in active engagement with a project of objective worth. For the proof to reach its conclusion, it must also be true that there is such a life among those available to the agent. The sentence does not say this; the proof assumes it.

\lem{7.1}{A life that includes engagement with $p$}

\paragraph{The hidden premise.} The phrase \enquote{a life that includes engagement with $p$} is used as if there were such a life available to the agent. Nothing in Definition~1 or Assumption~3 says that there is. Assumption~3 says that the agent is able to engage in $p$, which is a fact about her abilities; Theorem~4 says that a life of a certain kind is available to her, which is a fact about the lives she can lead. The step between them is a principle connecting abilities with available lives.

\begin{proposition}[The bridge]\label{p:bridge}
The proof of Theorem~4 requires the following principle, which is not stated in EXI-01-a: if an agent is able to engage in a project, then some life available to her includes engagement with that project.
\end{proposition}

\begin{proof}
Without the principle, Assumption~3 and Definition~1 are consistent with the falsity of Theorem~4. Suppose that an agent is able to engage in a worthwhile project, but that no life available to her includes engagement with it. Then no life available to her consists to any extent in engagement with that project, and if it is her only such project, Definition~1 gives no life available to her a positive degree of meaning. So the principle is required. It suffices with the other premises, since it supplies the life whose existence the second sentence of the proof assumes.
\end{proof}

The Lean module makes the principle an axiom, \texttt{engagement\_realizable}, and \texttt{\#print axioms} lists it among the axioms of the theorem (Section~\ref{s:lean}). The formalization thus exposes a premise that the prose leaves implicit. That is one of the services formalization performs, and in this case it shows where the modal content of \enquote{available} comes from.

\paragraph{Is the bridge true?} The principle connects two modalities: what an agent is able to do and what lives she can lead. It is true on any reading on which they are the same modality, since to be able to engage in a project is then to be able to lead a life that includes the engagement. It could fail only if the agent could engage in a project but could not lead a life that included it, which would require that abilities be exercisable only in lives the agent cannot lead. On the reading of \enquote{able} given at \Lm{4.4}, ability is what the agent can do from her actual situation, and what she can do from her situation is what some continuation of her life available to her includes. The principle is then analytic of the notions involved, and we think it is properly regarded as a clarification of \enquote{available} rather than as a substantive premise. Its formalization as an axiom is nonetheless correct, since the module does not define availability in terms of ability, and an undefined connection must be stated.

\paragraph{\enquote{engagement} and \enquote{active engagement}.} The phrase says \enquote{engagement with $p$}; the rest of the sentence says \enquote{active engagement}. The proof moves from a life that includes engagement to a life that consists in active engagement. The move is valid only if the engagement included is active. The text supplies no argument for this, and none is needed if the bridge principle is read as concerning active engagement, since Assumption~3's \enquote{engage}, as the comment on \Lm{4.4} noted, refers to the ability whose exercise is active engagement in the sense of Definition~1. A life that included only half-hearted or perfunctory engagement with $p$ would not be the life the proof needs, and the bridge principle should be read as supplying a life of active engagement. The Lean module's predicate \texttt{Includes} carries this reading in its docstring: a life includes \enquote{active engagement} with a project.

\lem{7.2}{consists, to some positive extent, in}

\paragraph{From inclusion to extent.} The inference from \enquote{includes} to \enquote{consists to some positive extent in} is the second inference that the sentence contains. It depends on the measure of Definition~1. On a reading on which momentary engagement could be included in a life without making up any positive extent of it, the inference would fail: a single moment of engagement might be a part of measure zero. Proposition~\ref{p:neutral} states the condition under which the inference holds: the engagement included must be sustained over a non-null part of the life. The word \enquote{project}, read in the broad sense, supplies this condition, since a project is sustained engagement (Proposition~\ref{p:broad}), and so does the reading of \enquote{active engagement} as caring activity with a temporal structure.

\paragraph{\enquote{to some positive extent}.} The phrase states only positivity, which is all that Theorem~4 needs. It is the textual counterpart of the minimal claim of \Lm{5.2}. It also shows that the proof is insensitive to the worth of $p$, so long as $p$ has some: a project of slight worth and a project of great worth both yield a life of positive meaning, though not of the same degree.

\lem{8}{By Definition 1, it is meaningful to a positive degree.}

The last sentence applies Definition~1 from right to left: since the life consists to a positive extent in active engagement with a project of objective worth, the degree $d$ that Definition~1 assigns it is positive. The pronoun \enquote{it} refers to the life of \Lm{7}. Since that life is available to the agent by the bridge principle, a life of positive meaning is available, which is Theorem~4. The square that closes the proof marks the end of the derivation and not of the argument: the derivation is complete, but its premises, as Remark~5 goes on to say, are where the philosophical questions lie.

\paragraph{The proof as a whole.} The proof is valid, given the bridge principle, and short, because the theorem says little more than its premises. That is not a defect. The text was written to make the structure of the availability claim visible, and the proof shows that structure: the claim holds if and only if, for every agent, some project of objective worth is within her reach (given the bridge principle and Definition~1). Every objection to the availability of meaning must therefore be an objection to Definition~1, to Assumption~3 or to the bridge, and the Commentary has examined each. The proof also shows that Assumption~2 plays no role in the derivation (Proposition~\ref{p:redundant}), so that an objection to the availability claim that attacks the existence of objective worth is, strictly, an objection to Assumption~3, which entails it.

\subsection{Other routes to Theorem 4}\label{s:routes}

The proof of EXI-01-a is one route to Theorem~4. The collection contains the materials for others, and comparing them shows which premises are essential.

\paragraph{The route of EXI-01-c.} The alternative proof replaces Assumption~3 by Reach and Lemma~2.3 of EXI-01-c, and derives Assumption~3 as a corollary. It keeps Definition~1, and it needs the bridge principle at the same point as the original proof. Its premises are not fewer than those of EXI-01-a, but its premise about agents is a claim about objects within reach rather than about projects, and the cases discussed in Section~\ref{s:hard} show that it is the better supported of the two.

\paragraph{The route through achievement.} One might try to derive Theorem~4 from EXI-17-a's Corollary~3.6, that meaning is not hostage to luck in outcomes, together with EXI-04-a's Lemma~3.1, that every agent capable of intentional action takes some ends seriously. The attempt fails, and the way it fails is instructive. Corollary~3.6 removes one obstacle to availability, the dependence of meaning on success, but it does not supply a worthwhile project for every agent. Lemma~3.1 of EXI-04-a supplies, for every agent, ends taken seriously, but not ends of worth. To reach Theorem~4, one needs in addition that every agent can take seriously some end of worth, which is Assumption~3 again, or Reach. The route through achievement therefore shows that no combination of results about engagement and success suffices without a premise about the reach of worth. EXI-17-a says as much when it states that its corollary supports the availability claim without by itself establishing it.

\paragraph{The route through the attitude toward one's condition.} The comment on \Lm{4.4} considered Frankl's third way: the attitude one takes toward unavoidable suffering. If bearing one's condition with courage is a project of objective worth, and every agent capable of intentional action is able to take some attitude toward her condition, then every agent is able to engage in a project of objective worth, and Theorem~4 follows. This route has the weakest premise about reach, since an attitude toward one's condition is available whenever agency is, and the strongest premise about worth, since it requires that the manner in which one bears one's condition be itself a project of objective worth. We think the premise about worth is true, since courage is a good, but the route is less secure than EXI-01-c's because that premise is more contested than the worth of persons and of understanding. It has the merit of covering the limiting cases, such as prolonged solitary confinement, in which EXI-01-c's first route is closed and its second is narrow.

\paragraph{What every route needs.} Every route to Theorem~4 needs three things: an account of meaning on which engagement with worth suffices for positive meaning (Definition~1, read right to left), a premise that every agent has something of worth within her reach (Assumption~3, Reach, or the premise of the third route), and the bridge from what an agent can do to the lives available to her. None needs Assumption~2 in its derivation. This is a structural fact about the availability claim, and the Lean module displays it for the route of EXI-01-a.
